% ===================================================================== Chapter 5
\chapter{CAD and Computational Model}\label{ch:cad}

\section{SpaceClaim Geometry}
The geometry was built in ANSYS SpaceClaim entirely by an IronPython script run in batch mode, so that re-running one file
reproduces it. The solid annulus is sketched as two concentric circles ($R$ = 20 and 10~mm) and extruded 600~mm. SpaceClaim resolves the two
circles into a single annular face, so the bore produces no body of its own; a first build therefore contained only one body. The
fluid body is created by a second sketch and extrusion of the $R$ = 10~mm disc. An assertion on the body count caught the behaviour.

\section{Fluid-Solid Domain Creation}
The model contains two coincident bodies: \texttt{SOLID\_DOMAIN} (4 faces) and \texttt{FLUID\_DOMAIN} (3 faces). The fluid is a real body, not a
void, so that it can be meshed and becomes a Fluent cell zone. \texttt{ShareTopology.FindAndFix} made the $\diameter$20 interface conformal,
which lets Fluent create a coupled wall/wall-shadow pair without interpolation. Figure~\ref{fig:cad} shows renders of the exported geometry.

\begin{figure}[htbp]
\centering
\begin{subfigure}[t]{0.32\linewidth}\centering\includegraphics[width=\linewidth]{01_full_cad_model.png}\caption{Full model}\end{subfigure}\hfill
\begin{subfigure}[t]{0.30\linewidth}\centering\includegraphics[width=\linewidth]{02b_cross_section_end_on.png}\caption{End-face cross-section}\end{subfigure}\hfill
\begin{subfigure}[t]{0.32\linewidth}\centering\includegraphics[width=\linewidth]{05_fluid_solid_interface.png}\caption{Fluid--solid interface}\end{subfigure}
\caption[CAD model]{CAD model: solid annulus enclosing the fluid body, and the shared $\diameter$20 interface.}
\label{fig:cad}
\src{\provcad{} The renders use the reference STL tessellation (18 facets per circle), which is why the circles appear polygonal; the meshing source is the native model.}
\end{figure}

\section{Named Selections}
Eleven named selections were created by classifying faces by area and bounding box (Table~\ref{tab:namedsel}). They carry every boundary
condition of the CFD and structural models. The structural support face is deliberately an end face only; restraining the cylindrical wall
would suppress the free growth that LC1 is meant to measure.

%% generated by gen_tables.py - RE-ANALYSIS 2026
\begin{table}[htbp]
\centering
\footnotesize
\caption{Named selections created in SpaceClaim and carried into Fluent and Mechanical}
\label{tab:namedsel}
\begin{tabular}{l>{\raggedright\arraybackslash}p{0.30\linewidth}>{\raggedright\arraybackslash}p{0.32\linewidth}}
\toprule
Name & Entity & Use \\
\midrule
\texttt{FLUID\_DOMAIN} & body: air passage & Fluent fluid cell zone \\
\texttt{SOLID\_DOMAIN} & body: Inconel 718 annulus & Fluent solid zone; Mechanical body \\
\texttt{FLUID\_INLET} & \O{}20 disc at $z = 0$ & velocity inlet \\
\texttt{FLUID\_OUTLET} & \O{}20 disc at $z = 600$ mm & pressure outlet \\
\texttt{FLUID\_WALL} & \O{}20 cylinder, fluid side & near-wall grading, $y^+$ monitoring \\
\texttt{FLUID\_SOLID\_INTERFACE} & both sides of the \O{}20 surface & conjugate coupling \\
\texttt{SOLID\_INNER\_INTERFACE} & \O{}20 cylinder, solid side & conjugate coupling \\
\texttt{HEATED\_OUTER\_WALL} & \O{}40 cylinder & heat flux 8000 W/m$^2$ \\
\texttt{SOLID\_INLET\_END} & annulus at $z = 0$ & adiabatic \\
\texttt{SOLID\_OUTLET\_END} & annulus at $z = 600$ mm & adiabatic \\
\texttt{STRUCTURAL\_SUPPORT} & annulus at $z = 0$ & structural restraint face \\
\bottomrule
\end{tabular}
\par\smallskip\parbox{\linewidth}{\footnotesize\raggedright Source: \texttt{03\_CAD\_Geometry/CAD\_NOTES.md} \S6. All 11 selections survived the transfer to Workbench Meshing without re-creation.}
\end{table}


\section{Geometry Verification}
Every dimension, area and volume of the saved model was compared with its analytical value (Table~\ref{tab:geomver}). All agree to
double-precision round-off (largest relative deviation $1.8\times10^{-14}$\,\%). The model has 7 faces, no open surfaces, no sliver faces and
no defeaturing needs.

%% generated by gen_tables.py - RE-ANALYSIS 2026
\begin{table}[htbp]
\centering
\footnotesize
\caption{Geometry verification: analytical values against the SpaceClaim model}
\label{tab:geomver}
\begin{tabular}{lrrl}
\toprule
Quantity & Analytical & CAD & Relative difference \\
\midrule
Fluid cross-sectional area [mm$^2$] & 314.1592654 & 314.1592654 & $1.8\times10^{-14}$\,\% \\
\texttt{FLUID\_DOMAIN} volume [m$^3$] & $1.884955592\times10^{-4}$ & $1.884955592\times10^{-4}$ & $1.4\times10^{-14}$\,\% \\
\texttt{SOLID\_DOMAIN} volume [m$^3$] & $5.654866776\times10^{-4}$ & $5.654866776\times10^{-4}$ & 0 \\
\texttt{HEATED\_OUTER\_WALL} area [m$^2$] & 0.07539822369 & 0.07539822369 & 0 \\
Interface area, fluid side = solid side [m$^2$] & 0.03769911184 & 0.03769911184 & 0 (identical to 17 digits) \\
Fluid + solid volume [m$^3$] & \multicolumn{3}{l}{$7.539822369\times10^{-4}$ = full \O{}40 $\times$ 600 mm cylinder (no gap, no overlap)} \\
\bottomrule
\end{tabular}
\par\smallskip\parbox{\linewidth}{\footnotesize\raggedright Source: \texttt{03\_CAD\_Geometry/CAD\_NOTES.md} \S5 and \S7; raw values in \texttt{Geometry\_Check/cad\_measurements.json}.}
\end{table}


\section{Volume Verification}
The strongest single check is the volume sum: fluid plus solid volume equals the volume of the full $\diameter$40 $\times$ 600~mm cylinder,
$7.539822369\times10^{-4}$ m$^3$, to machine precision. An overlap would exceed this value and a gap would fall short of it. The interface
areas on the fluid and solid sides are identical to 17 significant digits.

\section{Workbench Architecture}
The Workbench project links a Geometry system (the native SpaceClaim file, shared), a Mesh system and a Fluent system; it was built by a
Workbench journal in batch mode. All eleven named selections transferred automatically. The CFD meshes themselves were generated by a
Python script (Chapter~\ref{ch:mesh}) and written in the Fluent mesh format. For the structural phase, the Workbench project was saved under
a new name and extended with an External Data system (temperature transfer) and Static Structural systems for LC1 and LC2; linear buckling,
the structural mesh study and each parametric case used further save-as copies, so that every earlier project stayed byte-identical.

\section{Model Preparation}
Before any solve the geometry was checked for sweepability: both domains have constant cross-sections over 600~mm, so hexahedral sweeps apply
directly and no tetrahedra are needed. The parametric thickness cases T01 and T03 were built from a parameterised copy of the same script; their
dimensions, areas and volumes also match the analytical values to round-off. Parasolid export was unavailable in the Student installation;
the native file and STEP were used instead.
